% ============================================================================
% Fichier  : partie4/P04_C17_bia.tex
% Titre    : L'Analyse d'Impact sur l'Activité
% Auteur   : MINKA MI NGUIDJOÏ Thierry Emmanuel
% Version  : 1.0 -- novembre 2026
% ============================================================================
% Sources : ISO 22301:2019 (8.2.2) ; ISO/TS 22317:2021. Analyse d'impact de
% la BANQUE SANAGA (Opération WOURI) : cas fictif.
% ============================================================================

\chapter{L'Analyse d'Impact sur l'Activité}
\label{chap:bia}
\chaptermark{Analyse d'impact sur l'activité}

\epigraph{%
    « On ne reprend pas tout en même temps~;\\
    on reprend d'abord ce qui ne peut pas attendre. »%
}{Principe du cours}

\niveauChapitre{Intermédiaire}{Chapitres~\ref{chap:matrices}
et~\ref{chap:incidents}~; cours SEC3266, chapitre sur la continuité.}

\begin{boxobjectifs}
\begin{itemize}
    \item Conduire une analyse d'impact sur l'activité selon l'ISO/TS~22317.
    \item Mesurer l'évolution des impacts dans le temps.
    \item Fixer la période d'interruption maximale tolérable, le délai de
          reprise et la perte de données admissible.
    \item Recenser les dépendances d'une activité.
    \item En déduire l'ordre de reprise.
\end{itemize}
\end{boxobjectifs}

\bigskip

% ============================================================================
\section*{Question directrice : quelles activités de la BANQUE SANAGA ne peuvent pas attendre~?}
% ============================================================================

La BANQUE SANAGA exploite son système central à Douala. M.~EYENGA,
responsable de la continuité, doit dire à la direction, pour chaque
activité, combien de temps elle peut s'arrêter et combien de données elle
peut perdre, avant de choisir et de financer un site de secours. Sans cette
analyse, tout semble urgent, et rien n'est priorisé.

\medskip
\textbf{Plan du chapitre.}
\begin{enumerate}
    \item L'objet de l'analyse d'impact (section~\ref{sec:c17-objet}).
    \item Mesurer l'impact dans le temps (section~\ref{sec:c17-temps}).
    \item Les objectifs de reprise (section~\ref{sec:c17-objectifs}).
    \item Les dépendances (section~\ref{sec:c17-dependances}).
    \item L'ordre de reprise (section~\ref{sec:c17-ordre}).
\end{enumerate}

\separateur

% ============================================================================
\section{L'Objet de l'Analyse d'Impact}
\label{sec:c17-objet}
% ============================================================================

\begin{boxdef}
\termecle{Analyse d'impact sur l'activité} (BIA)~: processus qui analyse
l'impact d'une interruption sur l'organisation au fil du temps, pour fixer
les priorités et les objectifs de reprise de chaque
activité~\cite{iso22301_2019,iso22317_2021}.
\end{boxdef}

L'analyse d'impact précède le choix des stratégies~: on ne finance un site
chaud que pour les activités qui le justifient. Elle se conduit avec les
\emph{responsables d'activité}, et non avec l'informatique seule~: c'est le
métier qui connaît l'impact d'une heure d'arrêt.

% ============================================================================
\section{Mesurer l'Impact dans le Temps}
\label{sec:c17-temps}
% ============================================================================

Pour chaque activité, on cote l'impact d'une interruption après une heure,
quatre heures, un jour, trois jours, une semaine, sur plusieurs
critères~: financier, réglementaire, clients, réputation. L'impact retenu à
chaque durée est le plus élevé des critères.

\begin{tableaucours}{L{4.0cm}ccccc}{%
    \tableauHeader{Activité} &
    \tableauHeader{1~h} &
    \tableauHeader{4~h} &
    \tableauHeader{1~j} &
    \tableauHeader{3~j} &
    \tableauHeader{1~sem.}
}
Virements et compensation & 2 & 4 & 4 & 4 & 4 \\
Retraits aux distributeurs, cartes & 1 & 2 & 3 & 4 & 4 \\
Opérations aux guichets & 1 & 2 & 3 & 4 & 4 \\
Banque en ligne & 1 & 1 & 2 & 3 & 4 \\
Instruction des crédits & 1 & 1 & 1 & 2 & 3 \\
Rapports réglementaires & 1 & 1 & 1 & 1 & 3 \\
\end{tableaucours}
\vspace{-0.8em}
\begin{center}{\small\itshape Tableau~17.1. Impact dans le temps (1 mineur à 4 critique), BANQUE SANAGA}\end{center}

\textbf{Lecture.} Les virements deviennent critiques en quatre heures,
parce que la séance de compensation du jour serait manquée~; les rapports
réglementaires peuvent attendre plusieurs jours. L'impact croît par
paliers, liés à des échéances externes. La direction a fixé que l'impact
devient inacceptable lorsqu'il atteint 3 (majeur)~: la première durée où
c'est le cas donne la période d'interruption maximale tolérable.

% ============================================================================
\section{Les Objectifs de Reprise}
\label{sec:c17-objectifs}
% ============================================================================

\begin{boxdef}
\termecle{Période d'interruption maximale tolérable} (MTPD)~: durée au-delà
de laquelle l'impact d'une interruption devient inacceptable.
\termecle{Délai de reprise visé} (RTO)~: durée, inférieure à la MTPD, au
terme de laquelle l'activité doit avoir repris à un niveau minimal.
\termecle{Perte de données maximale admissible} (RPO)~: quantité de données,
exprimée en durée, que l'on accepte de perdre. Le \termecle{niveau minimal
de continuité} (MBCO) est le niveau d'activité acceptable pendant la
perturbation.
\end{boxdef}

\begin{tableaucours}{L{4.0cm}R{1.6cm}R{1.6cm}R{1.8cm}L{2.8cm}}{%
    \tableauHeader{Activité} &
    \tableauHeader{MTPD} &
    \tableauHeader{RTO} &
    \tableauHeader{RPO} &
    \tableauHeader{Niveau minimal}
}
Virements et compensation & 4~h & 2~h & 0 & Virements urgents, séance
du jour \\
Retraits, cartes & 1~j & 4~h & 15~min & Distributeurs des grandes villes \\
Guichets & 1~j & 8~h & 15~min & Retraits plafonnés, mode manuel \\
Banque en ligne & 3~j & 12~h & 1~h & Consultation des comptes \\
Instruction des crédits & 1~sem. & 3~j & 1~j & Dossiers urgents \\
Rapports réglementaires & 1~sem. & 5~j & 1~j & Rapports à échéance \\
\end{tableaucours}
\vspace{-0.8em}
\begin{center}{\small\itshape Tableau~17.2. Objectifs de reprise de la BANQUE SANAGA (Opération WOURI)}\end{center}

\begin{boxalert}
\textbf{Violation typique.} Fixer le RTO égal à la MTPD.

\smallskip
\textbf{Règle~:} le RTO laisse une marge sous la MTPD, parce que la
détection, la décision et la bascule prennent du temps. Un RTO de 4~h pour
une MTPD de 4~h signifie qu'au moindre retard, l'impact devient
inacceptable.
\end{boxalert}

% ============================================================================
\section{Les Dépendances}
\label{sec:c17-dependances}
% ============================================================================

Une activité ne reprend que si toutes les ressources dont elle dépend
reprennent~: personnes, locaux, applications, données, fournisseurs. Le
RTO d'une ressource doit être au plus égal au plus petit RTO des
activités qui en dépendent.

\begin{tableaucours}{L{3.6cm}L{4.6cm}L{2.0cm}L{2.0cm}}{%
    \tableauHeader{Ressource} &
    \tableauHeader{Activités qui en dépendent} &
    \tableauHeader{RTO exigé} &
    \tableauHeader{RTO actuel}
}
Système central & Toutes & 2~h & 24~h \\
Annuaire et authentification & Toutes & 2~h & 6~h \\
Réseau entre agences & Guichets, retraits & 4~h & 4~h \\
Prestataire monétique & Retraits, cartes & 4~h & 2~h \\
Connexion aux systèmes de paiement & Virements & 2~h & 8~h \\
\end{tableaucours}
\vspace{-0.8em}
\begin{center}{\small\itshape Tableau~17.3. Dépendances et écarts de RTO}\end{center}

\begin{boxscenario}{Opération WOURI. Les écarts}
\textbf{Constatation.} Trois ressources ont un RTO actuel supérieur au RTO
exigé~: le système central (24~h pour 2~h), l'annuaire (6~h pour 2~h) et
la connexion aux systèmes de paiement (8~h pour 2~h).

\medskip
\textbf{Analyse} (confiance élevée). Avec ces écarts, une perte du site de
Douala manquerait la séance de compensation du jour~: la MTPD des
virements serait dépassée. L'analyse d'impact justifie un site de secours
chaud pour ces trois ressources, et pas pour les autres
(Chapitre~\ref{chap:pca}).
\end{boxscenario}

% ============================================================================
\section{L'Ordre de Reprise}
\label{sec:c17-ordre}
% ============================================================================

\begin{boxprocedure}
\textbf{Établir l'ordre de reprise}
\begin{enumerate}
    \item Classer les activités par RTO croissant.
    \item Pour chacune, lister ses ressources et leurs dépendances.
    \item Reprendre d'abord les ressources dont dépendent les activités les
          plus urgentes (réseau, annuaire, système central).
    \item Puis les activités, dans l'ordre des RTO.
    \item Vérifier que la somme des durées de reprise tient dans les RTO.
\end{enumerate}
\end{boxprocedure}

\begin{boxobjet}{R-W1 · Perte du site de Douala}
\textbf{Scénario~:} du fait d'un site unique, un incendie ou un
rançongiciel rend indisponible le système central, ce qui arrête les
virements au-delà de la séance de compensation. \textbf{Cotation~:}
V=2, I=4 (score~8, mais impact critique en quatre heures).
\textbf{Traitement~:} site de secours chaud pour le système central,
l'annuaire et la connexion aux paiements. \textbf{Propriétaire~:} le
directeur des opérations bancaires. \textbf{Indicateur~:} durée de bascule
mesurée au dernier test.
\end{boxobjet}

\begin{boxmetier}{Responsable de la continuité}
\textbf{Faire signer les objectifs par les métiers.} M.~EYENGA fait valider
chaque MTPD, RTO et RPO par le responsable de l'activité et présente à la
direction le coût de chaque niveau~: c'est la direction qui choisit, en
connaissance du prix.
\end{boxmetier}

\separateur

\begin{boxresume}
\textbf{Ce qu'il faut retenir de ce chapitre~:}
\begin{itemize}
    \item \textbf{L'analyse d'impact} précède les stratégies et se conduit
          avec les métiers.
    \item \textbf{L'impact croît dans le temps}, par paliers liés aux
          échéances.
    \item \textbf{MTPD, RTO, RPO, niveau minimal}~: le RTO laisse une marge
          sous la MTPD.
    \item \textbf{Les dépendances} imposent leur RTO aux ressources.
    \item \textbf{L'ordre de reprise} commence par les ressources des
          activités les plus urgentes.
\end{itemize}
\end{boxresume}

\begin{boxinfo}
\textbf{Points de contrôle de l'élève.} Je sais~: définir l'analyse
d'impact~; coter l'impact dans le temps~; distinguer MTPD, RTO, RPO et
niveau minimal~; repérer un écart de RTO sur une ressource~; établir un
ordre de reprise.
\end{boxinfo}

\bigskip

\begin{boxexo}[title={\ding{45}\quad Exercice~17.1. Définitions \niveaubase}]
Expliquez la différence entre MTPD et RTO avec l'exemple des virements.
\end{boxexo}

\begin{boxexo}[title={\ding{45}\quad Exercice~17.2. Lire le tableau \niveaubase}]
D'après le Tableau~17.1, au bout de combien de temps les retraits
deviennent-ils critiques~?
\end{boxexo}

\begin{boxexo}[title={\ding{45}\quad Exercice~17.3. RPO nul \niveaubase}]
Que signifie un RPO nul pour les virements~? Quelle technique permet de
l'atteindre~?
\end{boxexo}

\begin{boxexo}[title={\ding{45}\quad Exercice~17.4. Écarts \niveaumoyen}]
Calculez, pour chaque ressource du Tableau~17.3, l'écart entre RTO actuel et
RTO exigé. Laquelle est la plus urgente~?
\end{boxexo}

\begin{boxexo}[title={\ding{45}\quad Exercice~17.5. Le prestataire monétique \niveaumoyen}]
Le prestataire monétique a un RTO de 2~h, meilleur que l'exigence. Cela
suffit-il pour que les retraits reprennent en 4~h~? Justifiez par les
dépendances.
\end{boxexo}

\begin{boxexo}[title={\ding{45}\quad Exercice~17.6. L'ordre \niveaumoyen}]
Établissez l'ordre de reprise des ressources et des activités de la BANQUE
SANAGA.
\end{boxexo}

\begin{boxexo}[title={\ding{45}\quad Exercice~17.7. Le niveau minimal \niveaumoyen}]
Proposez un niveau minimal de continuité pour les guichets pendant trois
jours sans système central.
\end{boxexo}

\begin{boxexo}[title={\ding{45}\quad Exercice~17.8. LiZbiZ \niveaumoyen}]
Construisez le tableau de l'impact dans le temps pour la facturation et le
suivi des expéditions de LiZbiZ.
\end{boxexo}

\begin{boxexo}[title={\ding{45}\quad Exercice~17.9. Le questionnaire \niveauexpert}]
Rédigez le questionnaire d'entretien d'analyse d'impact que M.~EYENGA
soumet aux responsables d'activité.
\end{boxexo}

\begin{boxexo}[title={\ding{45}\quad Exercice~17.10. Le coût des objectifs \niveauexpert}]
Présentez à la direction de la BANQUE SANAGA trois options de RTO pour le
système central (2~h, 8~h, 24~h), avec leurs conséquences sur les
activités et un ordre de grandeur de coût.
\end{boxexo}

\bigskip

\begin{boxinfo}
\textbf{Pour aller plus loin.}
\begin{itemize}
    \item ISO~22301~\cite{iso22301_2019}~; ISO/TS~22317~\cite{iso22317_2021}.
\end{itemize}
\end{boxinfo}

\begin{boxinfo}
\textbf{Transition.} Les objectifs fixés, il faut choisir les stratégies et
écrire les plans. C'est l'objet du Chapitre~\ref{chap:pca}.
\end{boxinfo}
